\documentclass[11pt,a4paper]{article}

\usepackage[T1]{fontenc}
\usepackage[utf8]{inputenc}
\usepackage{lmodern}

\usepackage[a4paper,margin=1in]{geometry}
\usepackage{setspace}
\setstretch{1.12}

\setlength{\parindent}{15pt}
\setlength{\parskip}{0pt}

\usepackage{lmodern}
\usepackage{amsmath,amssymb,amsfonts}
\usepackage{booktabs}
\usepackage{longtable}
\usepackage{array}
\usepackage{tabularx}
\usepackage{enumitem}
\usepackage{hyperref}
\usepackage{setspace}
\usepackage{ragged2e}

\hypersetup{
	colorlinks=true,
	linkcolor=black,
	citecolor=black,
	urlcolor=blue,
	pdftitle={External Verification Checklist for the General Theory of Cognitive Structuring},
	pdfauthor={Kostiantyn Osmolovskyi}
}

\setlist[itemize]{topsep=3pt,itemsep=2pt,parsep=0pt}
\setlist[enumerate]{topsep=3pt,itemsep=2pt,parsep=0pt}


\title{External Verification Checklist\\ for the General Theory of Cognitive Structuring}
	
\author{
	Kostiantyn Osmolovskyi\\
	\small Independent Researcher, Ukraine\\
	\small ORCID: 0009-0006-3144-7237\\
	\small \texttt{constantinosmol@gmail.com}
}

\date{}

% ---------------------------------------------------------------------------
% Part of a series: General Theory of Cognitive Structuring (GTCS)
% Autor: Kostiantyn Osmolovskyi (ORCID: 0009-0006-3144-7237)
% License: Creative Commons Attribution 4.0 International (CC BY 4.0)
% Repository: https://github.com/k-osmolovskyi/k-osmolovskyi.github.io
% DOI: https://doi.org/10.5281/zenodo.19701916
% ---------------------------------------------------------------------------

\begin{document}
	\maketitle
	
\begin{center}
	\small
	Technical Note No. 2026-7\\
	Version 1.0
\end{center}

\vspace{0.5em}	
	\section{Purpose of this checklist}
	
	This checklist is intended as a practical aid for external readers evaluating the
	\emph{General Theory of Cognitive Structuring}. It is designed for reviewers, editors,
	formal verifiers, adjacent-field researchers, and other readers who want a compact
	verification instrument rather than a full interpretive or technical companion.
	
	The checklist does not replace the primary papers or the companion documents.
	Its function is narrower:
	\begin{itemize}[leftmargin=2em]
		\item to identify the main verification questions,
		\item to reduce accidental omissions during external review,
		\item to separate different kinds of verification tasks,
		\item to provide a compact record of what has and has not been checked.
	\end{itemize}
	
	\section{How to use this checklist}
	
	The checklist may be used in at least three ways:
	\begin{enumerate}[leftmargin=2em]
		\item as a quick review instrument after reading the companion documents;
		\item as a structured note-taking form during external verification;
		\item as a compact audit sheet for collaborative or repeated checking.
	\end{enumerate}
	
	The document is organized by verification domain rather than by paper order. A reader may
	complete only the sections relevant to the question under review.
	
	\section{Response scale}
	
	The following response scale may be used throughout the checklist:
	
	\begin{itemize}[leftmargin=2em]
		\item \textbf{Yes} --- clearly satisfied;
		\item \textbf{Mostly} --- substantially satisfied, with minor ambiguity;
		\item \textbf{Unclear} --- not yet clear from the current materials;
		\item \textbf{No} --- not satisfied or contradicted;
		\item \textbf{N/A} --- not applicable to the current verification task.
	\end{itemize}
	
	\vspace{0.5em}
	\noindent
	A brief comments field is included after each section.
	
	\section{Section A: Entry and scope verification}
	
	\begin{longtable}{>{\RaggedRight\arraybackslash}p{0.08\textwidth}
			>{\RaggedRight\arraybackslash}p{0.66\textwidth}
			>{\RaggedRight\arraybackslash}p{0.16\textwidth}}
		\toprule
		\textbf{ID} & \textbf{Verification item} & \textbf{Response} \\
		\midrule
		\endfirsthead
		\toprule
		\textbf{ID} & \textbf{Verification item} & \textbf{Response} \\
		\midrule
		\endhead
		
		A1 & Is the distinction between processing and structural updating stated clearly enough to function as an entry distinction? & \\
		
		A2 & Is the entry distinction narrower than ``cognitive change in general'' and explicitly limited to the availability of structural updating? & \\
		
		A3 & Does the entry paper avoid pretending to be a full mechanistic theory of structural change? & \\
		
		A4 & Is the role of admissibility introduced as a clarificatory condition rather than as a catch-all explanation of all non-change? & \\
		
		A5 & Are the limits of the entry paper made explicit, rather than left implicit? & \\
		
		\bottomrule
	\end{longtable}
	
	\textbf{Comments for Section A:}
	
	\vspace{3em}
	
	\section{Section B: Terminology and interpretive stability}
	
	\begin{longtable}{>{\RaggedRight\arraybackslash}p{0.08\textwidth}
			>{\RaggedRight\arraybackslash}p{0.66\textwidth}
			>{\RaggedRight\arraybackslash}p{0.16\textwidth}}
		\toprule
		\textbf{ID} & \textbf{Verification item} & \textbf{Response} \\
		\midrule
		\endfirsthead
		\toprule
		\textbf{ID} & \textbf{Verification item} & \textbf{Response} \\
		\midrule
		\endhead
		
		B1 & Are the main terms used in an operational, architectural, or regulatory sense rather than in an ordinary-language sense? & \\
		
		B2 & Is \emph{admissibility} kept distinct from truth, correctness, awareness, or mere possibility? & \\
		
		B3 & Is \emph{coherence} kept distinct from semantic consistency, narrative unity, or subjective harmony? & \\
		
		B4 & Is \emph{coherence representation} kept distinct from geometric coherence itself? & \\
		
		B5 & Is \emph{pre-symbolic admissibility} treated as a filtering or gating condition rather than as a content type? & \\
		
		B6 & Is \emph{manifestation} treated as expression of organization rather than the organization itself? & \\
		
		B7 & Is \emph{overload memory} kept distinct from autobiographical or semantic memory? & \\
		
		B8 & Are \emph{invariants} treated as structural constraints rather than automatically as explicit beliefs or verbal rules? & \\
		
		B9 & Are \emph{mode} and \emph{regime} used in controlled architectural ways rather than as loose state labels? & \\
		
		B10 & Are the glossary distinctions respected by the paper corpus rather than contradicted by it? & \\
		
		\bottomrule
	\end{longtable}
	
	\textbf{Comments for Section B:}
	
	\vspace{3em}
	
	\section{Section C: Separation of levels}
	
	\begin{longtable}{>{\RaggedRight\arraybackslash}p{0.08\textwidth}
			>{\RaggedRight\arraybackslash}p{0.66\textwidth}
			>{\RaggedRight\arraybackslash}p{0.16\textwidth}}
		\toprule
		\textbf{ID} & \textbf{Verification item} & \textbf{Response} \\
		\midrule
		\endfirsthead
		\toprule
		\textbf{ID} & \textbf{Verification item} & \textbf{Response} \\
		\midrule
		\endhead
		
		C1 & Is global structural condition kept distinct from regulatorily accessible condition? & \\
		
		C2 & Is accessibility kept distinct from representation? & \\
		
		C3 & Is representation kept distinct from manifestation? & \\
		
		C4 & Is manifestation kept distinct from underlying architecture? & \\
		
		C5 & Is geometric coherence kept distinct from local structural pressure signals? & \\
		
		C6 & Is structural admissibility of updating kept distinct from pre-symbolic admissibility of discrepancy? & \\
		
		C7 & Is overload kept distinct from overload memory? & \\
		
		C8 & Is structural updating kept distinct from learning, local adjustment, and processing within an existing organization? & \\
		
		C9 & Is inter-system incompatibility treated at the level claimed by the theory rather than collapsed into mere disagreement of representations? & \\
		
		\bottomrule
	\end{longtable}
	
	\textbf{Comments for Section C:}
	
	\vspace{3em}
	
	\section{Section D: Formal dependency and acyclicity}
	
	\begin{longtable}{>{\RaggedRight\arraybackslash}p{0.08\textwidth}
			>{\RaggedRight\arraybackslash}p{0.66\textwidth}
			>{\RaggedRight\arraybackslash}p{0.16\textwidth}}
		\toprule
		\textbf{ID} & \textbf{Verification item} & \textbf{Response} \\
		\midrule
		\endfirsthead
		\toprule
		\textbf{ID} & \textbf{Verification item} & \textbf{Response} \\
		\midrule
		\endhead
		
		D1 & Are the dependency criteria of the series stated explicitly enough to evaluate circularity? & \\
		
		D2 & Is logical dependency kept distinct from textual cross-reference? & \\
		
		D3 & Is retrospective clarification kept distinct from reverse logical generation? & \\
		
		D4 & Is synthetic integration kept distinct from foundational dependence? & \\
		
		D5 & Does the paper-level unfolding remain representable as a directed acyclic structure under the stated criteria? & \\
		
		D6 & Are apparent cycles explainable by refinement, synthesis, or later formal grounding rather than by actual circular dependence? & \\
		
		D7 & Does the \emph{General Theory of Cognitive Structuring} function as a convergence node rather than as the original source of the earlier layers? & \\
		
		D8 & Are node and edge types controlled well enough to prevent false edges? & \\
		
		\bottomrule
	\end{longtable}
	
	\textbf{Comments for Section D:}
	
	\vspace{3em}
	
	\section{Section E: Notation and formal traceability}
	
	\begin{longtable}{>{\RaggedRight\arraybackslash}p{0.08\textwidth}
			>{\RaggedRight\arraybackslash}p{0.66\textwidth}
			>{\RaggedRight\arraybackslash}p{0.16\textwidth}}
		\toprule
		\textbf{ID} & \textbf{Verification item} & \textbf{Response} \\
		\midrule
		\endfirsthead
		\toprule
		\textbf{ID} & \textbf{Verification item} & \textbf{Response} \\
		\midrule
		\endhead
		
		E1 & Is notation sufficiently stabilized across papers to permit cross-paper reading? & \\
		
		E2 & Are early and later versions of the stability region ($U$ vs.\ $U_X(I_t)$) readable as a controlled strengthening rather than as incompatible notions? & \\
		
		E3 & Are symbols for geometric coherence and coordination cost distinguishable by index and context? & \\
		
		E4 & Are representational symbols such as $\widehat{C}_t$, $Cb_t$, and related variables kept distinct from geometric coherence? & \\
		
		E5 & Are assumptions, definitions, operators, constructions, and derived results kept distinct in the technical appendix? & \\
		
		E6 & Can the direct prerequisites of the main definitions and operators be located without excessive reconstruction? & \\
		
		E7 & Can the main proposition- and theorem-level results be traced back to explicit prior objects and assumptions? & \\
		
		\bottomrule
	\end{longtable}
	
	\textbf{Comments for Section E:}
	
	\vspace{3em}
	
	\section{Section F: Assumptions and scope control}
	
	\begin{longtable}{>{\RaggedRight\arraybackslash}p{0.08\textwidth}
			>{\RaggedRight\arraybackslash}p{0.66\textwidth}
			>{\RaggedRight\arraybackslash}p{0.16\textwidth}}
		\toprule
		\textbf{ID} & \textbf{Verification item} & \textbf{Response} \\
		\midrule
		\endfirsthead
		\toprule
		\textbf{ID} & \textbf{Verification item} & \textbf{Response} \\
		\midrule
		\endhead
		
		F1 & Are global architectural assumptions distinguishable from local theorem-specific assumptions? & \\
		
		F2 & Are regularity assumptions stated where theorem-level results require them? & \\
		
		F3 & Are order-preserving, monotonicity, boundedness, compactness, or drift assumptions introduced only where needed? & \\
		
		F4 & Does the theory avoid silently upgrading assumption-dependent results into unconditional architectural claims? & \\
		
		F5 & Are the limits of each paper's contribution stated clearly enough to prevent over-reading? & \\
		
		F6 & Do later papers avoid retroactively overstating the scope of earlier ones? & \\
		
		\bottomrule
	\end{longtable}
	
	\textbf{Comments for Section F:}
	
	\vspace{3em}
	
	\section{Section G: Core architectural branch}
	
	\begin{longtable}{>{\RaggedRight\arraybackslash}p{0.08\textwidth}
			>{\RaggedRight\arraybackslash}p{0.66\textwidth}
			>{\RaggedRight\arraybackslash}p{0.16\textwidth}}
		\toprule
		\textbf{ID} & \textbf{Verification item} & \textbf{Response} \\
		\midrule
		\endfirsthead
		\toprule
		\textbf{ID} & \textbf{Verification item} & \textbf{Response} \\
		\midrule
		\endhead
		
		G1 & Is coherence defined geometrically relative to an admissible stability region? & \\
		
		G2 & Is structural admissibility defined over the joint regulatory state rather than over instantaneous tension alone? & \\
		
		G3 & Is overload memory introduced as a genuine state variable rather than as a restatement of present tension? & \\
		
		G4 & Is trajectory dependence treated as a structural consequence of overload memory rather than as a merely narrative claim about history? & \\
		
		G5 & Are invariants introduced as historically constituted structural constraints rather than as primitive semantic objects? & \\
		
		G6 & Does compression function as a constrained admissible transformation over invariant architecture? & \\
		
		G7 & Is architectural level defined geometrically rather than through content amount or intuitive depth? & \\
		
		G8 & Is level transition tied to non-equivalent change in admissible geometry rather than to mere increase of intensity or complexity? & \\
		
		\bottomrule
	\end{longtable}
	
	\textbf{Comments for Section G:}
	
	\vspace{3em}
	
	\section{Section H: Access and representation branch}
	
	\begin{longtable}{>{\RaggedRight\arraybackslash}p{0.08\textwidth}
			>{\RaggedRight\arraybackslash}p{0.66\textwidth}
			>{\RaggedRight\arraybackslash}p{0.16\textwidth}}
		\toprule
		\textbf{ID} & \textbf{Verification item} & \textbf{Response} \\
		\midrule
		\endfirsthead
		\toprule
		\textbf{ID} & \textbf{Verification item} & \textbf{Response} \\
		\midrule
		\endhead
		
		H1 & Is coherence representation introduced as an internal regulatory variable rather than as a reconstruction of full geometric coherence? & \\
		
		H2 & Is partial observability of structural stability used explicitly to motivate representation? & \\
		
		H3 & Is pre-symbolic admissibility introduced as an earlier filtering layer before representation? & \\
		
		H4 & Does the admissible discrepancy domain function as the domain from which regulatory accessibility is later derived? & \\
		
		H5 & Is restricted accessibility formulated as a consequence of geometric coherence plus admissibility-constrained access? & \\
		
		H6 & Is non-injectivity of regulatory accessibility treated as a structural consequence of admissibility constraints rather than as noise or approximation error? & \\
		
		H7 & Are representation, discrepancy filtering, and structural admissibility kept as distinct operators or levels? & \\
		
		\bottomrule
	\end{longtable}
	
	\textbf{Comments for Section H:}
	
	\vspace{3em}
	
	\section{Section I: Identity and multi-system branch}
	
	\begin{longtable}{>{\RaggedRight\arraybackslash}p{0.08\textwidth}
			>{\RaggedRight\arraybackslash}p{0.66\textwidth}
			>{\RaggedRight\arraybackslash}p{0.16\textwidth}}
		\toprule
		\textbf{ID} & \textbf{Verification item} & \textbf{Response} \\
		\midrule
		\endfirsthead
		\toprule
		\textbf{ID} & \textbf{Verification item} & \textbf{Response} \\
		\midrule
		\endhead
		
		I1 & Is identity derived as a structural consequence of bounded regulation rather than introduced as a primitive representational object? & \\
		
		I2 & Are identity-cores defined through low-overload or attractor-like regulatory structure rather than through narrative self-description? & \\
		
		I3 & Is multi-system comparability formulated without assuming shared metric coherence or shared configuration space? & \\
		
		I4 & Is inter-system order alignment treated as a minimal transferable structure rather than as representational identity? & \\
		
		I5 & Is inter-system conflict defined over admissible discrepancy structures rather than over representation mismatch alone? & \\
		
		I6 & Does the theory preserve the distinction between alignment and conflict as different extension problems? & \\
		
		\bottomrule
	\end{longtable}
	
	\textbf{Comments for Section I:}
	
	\vspace{3em}
	
	\section{Section J: Package usability}
	
	\begin{longtable}{>{\RaggedRight\arraybackslash}p{0.08\textwidth}
			>{\RaggedRight\arraybackslash}p{0.66\textwidth}
			>{\RaggedRight\arraybackslash}p{0.16\textwidth}}
		\toprule
		\textbf{ID} & \textbf{Verification item} & \textbf{Response} \\
		\midrule
		\endfirsthead
		\toprule
		\textbf{ID} & \textbf{Verification item} & \textbf{Response} \\
		\midrule
		\endhead
		
		J1 & Does the verification package reduce the amount of reconstruction required from an external reader? & \\
		
		J2 & Is there a clear reading order for the package? & \\
		
		J3 & Does each companion document have a distinct and non-redundant function? & \\
		
		J4 & Can a reviewer with limited time still verify the main architecture without reading every paper in full? & \\
		
		J5 & Is the package sufficient to support an informed external check of terminology, dependency structure, and formal traceability? & \\
		
		\bottomrule
	\end{longtable}
	
	\textbf{Comments for Section J:}
	
	\vspace{3em}
	
	\section{Summary sheet}
	
	The following compact summary may be used after completion of the checklist.
	
	\begin{longtable}{>{\RaggedRight\arraybackslash}p{0.42\textwidth}
			>{\RaggedRight\arraybackslash}p{0.18\textwidth}
			>{\RaggedRight\arraybackslash}p{0.30\textwidth}}
		\toprule
		\textbf{Verification domain} & \textbf{Overall assessment} & \textbf{Notes} \\
		\midrule
		\endfirsthead
		\toprule
		\textbf{Verification domain} & \textbf{Overall assessment} & \textbf{Notes} \\
		\midrule
		\endhead
		
		Entry and scope & & \\
		
		Terminology and interpretive stability & & \\
		
		Separation of levels & & \\
		
		Formal dependency and acyclicity & & \\
		
		Notation and formal traceability & & \\
		
		Assumptions and scope control & & \\
		
		Core architectural branch & & \\
		
		Access and representation branch & & \\
		
		Identity and multi-system branch & & \\
		
		Package usability & & \\
		
		\bottomrule
	\end{longtable}
	
	\section{Closing note}
	
	This checklist is intended to make external verification easier and more systematic, not more burdensome. It should be used as a compact instrument for checking whether the main distinctions, dependencies, and formal relations of the \emph{General Theory of Cognitive Structuring} remain visible and controlled for a reader approaching the series from outside its immediate development.


\end{document}